\documentclass[11pt]{article}
\usepackage{amsmath,amssymb,amsthm}
\newtheorem{theorem}{Theorem}
\newtheorem{lemma}[theorem]{Lemma}
\newtheorem{corollary}[theorem]{Corollary}
\newtheorem{conjecture}[theorem]{Conjecture}
\newtheorem{remark}[theorem]{Remark}
\newcommand{\EMD}{\operatorname{EMD}}
\newcommand{\rk}{\operatorname{rank}}

\title{The Master Recursion for $Q_{n,c}(q)$, the Preimage EMD Dichotomy,\\
and the Shape of the Level-$2$ Numerator}
\author{Seed 7, Layer 2, Round 2}
\date{2026-07-04}

\begin{document}
\maketitle

\begin{abstract}
Composing the (proved) adjugate inversion of the Corteel--Welsh system with the
conversion $Q_{n,c}=(q^\ell;q^\ell)_n\,g_{c,n}$ yields an explicit one-step
recursion expressing $Q_n$ in terms of $Q_{n-1}$ and $Q_{n-2}$ via
Earth-Mover's-Distance weights on profile space (Theorem~\ref{thm:master}).
We prove a dichotomy for the behaviour of $\EMD(c,\cdot)$ under the three
$2$-shift preimage moves (Lemma~\ref{lem:dichotomy}) and deduce that all
negativity in the level-$2$ numerator $N_2=(1+q^2+q^4)Q_2$ is confined to three
rigid ``sandwich'' shapes with a single negative coefficient each
(Theorem~\ref{thm:shape}). We also show that divisibility of $N_n$ by
$1+q^n+q^{2n}$ factors through a single scalar identity per root, generalizing
the level-$1$ EMD Equidistribution Theorem. Positivity of $N_n$ (verified for
$n\le 5$ at $d=4$ and $n\le 2$ at $d\in\{5,7\}$) remains conjectural.
\end{abstract}

\section{Setup and notation}

Fix $r=3$ and a profile $c=(c_0,c_1,c_2)$ with $d=c_0+c_1+c_2$,
$\ell=\gcd(d,3)$. Write $G_c(z,q)=(zq;q)_\infty F_c(z,q)$,
$g_{c,n}=[z^n]G_c(z,q)$ and $Q_{n,c}=(q^\ell;q^\ell)_n\,g_{c,n}$.
For $J\subseteq I_c=\{i:c_i>0\}$ let $c(J)$ be the Corteel--Welsh shifted
profile. The Earth Mover's Distance on profiles satisfies (proved in Round~2
Layer~1, Seed~4/Seed~7)
\begin{equation}\label{eq:emd}
\EMD(c,c') \;=\; 3\max\bigl(0,\;c'_1-c_1,\;c_0-c'_0\bigr)+(c'_0-c_0)-(c'_1-c_1),
\end{equation}
$\operatorname{adj}(I-A(x))[c,c']=x^{\EMD(c,c')}$ and
$\det(I-A(x))=-(x^3-1)$, whence the \emph{adjugate inversion}
\begin{equation}\label{eq:inv}
g_{c,n}\;=\;\frac{1}{1-q^{3n}}\sum_{c'} q^{\,n\,\EMD(c,c')}\,b_n(c').
\end{equation}

\section{The master recursion}

\begin{theorem}[Master recursion]\label{thm:master}
For $n\ge 1$ (with $Q_{-1}:=0$, $Q_{0,c}:=1$),
\[
Q_{n,c}\;=\;\frac{1-q^{\ell n}}{1-q^{3n}}
\sum_{c'} q^{\,n\,\EMD(c,c')}\,\mathrm{BR}_n(c'),
\]
where
\begin{align*}
\mathrm{BR}_n(c') \;=\;& q^{2n-1}\!\!\sum_{\substack{J\subseteq I_{c'}\\ |J|=2}}\!\! Q_{n-1,c'(J)}
\;+\;[\rk(c')=3]\Bigl(-(q^{3n-2}+q^{3n-1})\,Q_{n-1,c'}\\
&\qquad\qquad +\;q^{3n-3}\,(1-q^{\ell(n-1)})\,Q_{n-2,c'}\Bigr).
\end{align*}
For $\ell=1$ the prefactor is $1/(1+q^n+q^{2n})$.
\end{theorem}

\begin{proof}
Extract $[z^n]$ from the Corteel--Welsh equation for $G_c$,
$G_c(z,q)=\sum_{\emptyset\ne J\subseteq I_c}(-1)^{|J|-1}(zq;q)_{|J|-1}
G_{c(J)}(zq^{|J|},q)$, using $(zq;q)_0=1$, $(zq;q)_1=1-zq$,
$(zq;q)_2=1-z(q+q^2)+z^2q^3$. The $z^0$-coefficients of the truncated
Pochhammers produce the matrix $A(q^n)$ acting on the current level; the
remaining terms give
\[
b_n(c)\;=\;q^{2n-1}\!\!\sum_{\substack{J\subseteq I_c\\|J|=2}}\!\! g_{c(J),n-1}
\;-\;[\rk(c)=3]\,(q^{3n-2}+q^{3n-1})\,g_{c,n-1}
\;+\;[\rk(c)=3]\,q^{3n-3}\,g_{c,n-2},
\]
where we used that $|J|=3$ forces $c(J)=c$ and requires $\rk(c)=3$.
Insert this into \eqref{eq:inv} and substitute
$g_{\cdot,m}=Q_{m,\cdot}/(q^\ell;q^\ell)_m$, multiplying through by
$(q^\ell;q^\ell)_n$ and using
$(q^\ell;q^\ell)_n/(q^\ell;q^\ell)_{n-1}=1-q^{\ell n}$ and
$(q^\ell;q^\ell)_n/(q^\ell;q^\ell)_{n-2}=(1-q^{\ell n})(1-q^{\ell(n-1)})$.
At $n=1$ the $g_{c,-1}$ term is absent.
\end{proof}

\begin{remark}
Verified against an independent Neumann-iteration solve of the CW linear
system (no adjugate used): $d=4$, $n\le3$; $d=5,7$, $n\le2$; all profiles,
exact coefficient match, and $Q_{2,c}(1)=\bigl(\tfrac{(d+1)(d+2)}{6}-1\bigr)^2$.
The fully unfolded $n=2$ double sum, obtained by inserting the closed form
$Q_{1,c}=\frac{1}{1+q+q^2}\sum_{c'}q^{\EMD(c,c')}B(c')$
($B=q(2-q),\,q,\,0$ for ranks $3,2,1$), also matches exactly.
\end{remark}

\section{The preimage EMD dichotomy}

For a profile $c'$ call $c''$ a \emph{$2$-shift preimage} of $c'$ if
$c''(J)=c'$ for some $J\subseteq I_{c''}$ with $|J|=2$. Explicitly the three
candidate moves are
\[
J=\{0,1\}:\; c''=c'+e_0-e_2,\qquad
J=\{1,2\}:\; c''=c'-e_0+e_1,\qquad
J=\{0,2\}:\; c''=c'-e_1+e_2 .
\]
A direct support check shows: a rank-$3$ profile $c'$ has exactly three
preimages, a rank-$2$ profile exactly one, a rank-$1$ profile none.

\begin{lemma}[Preimage EMD dichotomy]\label{lem:dichotomy}
Fix $c$ and $c'$, put $e_0=\EMD(c,c')$, and for each existing preimage $c''$
put $\Delta=\EMD(c,c'')-e_0$. Then $\Delta\in\{-2,+1\}$, and at most one of
the three preimages has $\Delta=-2$.
\end{lemma}

\begin{proof}
Set $u=c'_0-c_0$, $v=c'_1-c_1$, so by \eqref{eq:emd}
$e_0=f(u,v):=3M(u,v)+u-v$ with $M(u,v)=\max(0,v,-u)$. The three preimage
moves act on $(u,v)$ by $a=(+1,0)$, $b=(-1,+1)$, $c=(0,-1)$.

\emph{Move $a$:} the linear part $u-v$ increases by $1$;
$M(u+1,v)=\max(0,v,-u-1)\in\{M,M-1\}$. Hence $\Delta_a\in\{1,-2\}$, and
$\Delta_a=-2$ iff $M=-u$ and $-u-1\ge\max(0,v)$, i.e.\
$-u\ge\max(0,v)+1$ (in particular $u\le-1$).

\emph{Move $c$:} the linear part increases by $1$;
$M(u,v-1)\in\{M,M-1\}$, and $\Delta_c=-2$ iff $v\ge\max(0,-u)+1$
(in particular $v\ge1$).

\emph{Move $b$:} the linear part decreases by $2$;
$M(u-1,v+1)=\max(0,v+1,-u+1)\in\{M,M+1\}$, so $\Delta_b\in\{-2,1\}$, and
$\Delta_b=-2$ iff $M(u-1,v+1)=M$, which forces $M=0$ with $v+1\le0$ and
$-u+1\le0$, i.e.\ $v\le-1$ and $u\ge1$.

Mutual exclusivity of the three conditions: $\Delta_a=-2$ needs $u\le-1$
while $\Delta_b=-2$ needs $u\ge1$; $\Delta_c=-2$ needs $v\ge1$ while
$\Delta_b=-2$ needs $v\le-1$; and $\Delta_a=-2$ with $\Delta_c=-2$ gives
$-u\ge v+1$ and $v\ge -u+1$, whence $-u\ge -u+2$, absurd.
\end{proof}

(Computationally re-checked for $d\in\{4,5,7,8,10\}$, all pairs $(c,c')$:
zero violations; both values $\{-2,1\}$ occur for ranks $2$ and $3$.)

\section{The shape of the level-2 numerator}

Let $\ell=1$ and define $N_2(c):=(1+q^2+q^4)\,Q_{2,c}
=\sum_{c'}q^{2\EMD(c,c')}\mathrm{BR}_2(c')$.
Regrouping the sum by the profile appearing as the argument of $Q_1$ and
applying Lemma~\ref{lem:dichotomy} gives:

\begin{theorem}[Shape theorem]\label{thm:shape}
Fix $c$ and write $e_0=e_0(c')=\EMD(c,c')$, and for rank-$3$ profiles let
$k(c,c')\in\{0,1\}$ be the number of preimages of $c'$ with $\Delta=-2$. Then
\begin{align*}
N_2(c)\;=\;&\sum_{\rk(c')=2} q^{2\EMD(c,c'')+3}\,Q_{1,c'}
\qquad (c''=\text{the unique preimage of }c')\\
&+\sum_{\substack{\rk(c')=3\\ k=0}} q^{2e_0+4}\,(2q-1)\,Q_{1,c'}
\;+\sum_{\substack{\rk(c')=3\\ k=1}} q^{2e_0-1}\,(1-q^5+q^6)\,Q_{1,c'}\\
&+\sum_{\rk(c')=3} q^{2e_0+3}\,(1-q).
\end{align*}
\end{theorem}

\begin{proof}
By Theorem~\ref{thm:master} at $n=2$,
$N_2(c)=\sum_{c'}W(c,c')\,Q_{1,c'}+\sum_{\rk(c')=3}q^{2e_0+3}(1-q)$ with
$W(c,c')=q^3\sum_{c''}q^{2\EMD(c,c'')}-[\rk(c')=3](q^4+q^5)q^{2e_0}$, the
inner sum over preimages $c''$ of $c'$; the last term is the $Q_0$
contribution $q^3(1-q)Q_{0,c'}=q^3(1-q)$. For rank-$2$ profiles $W$ is the
single monomial shown. For rank-$3$ profiles all three preimages exist and by
Lemma~\ref{lem:dichotomy} their exponents are $2(e_0+\Delta_i)+3$ with
$\Delta_i\in\{-2,1\}$ and at most one $-2$:
if $k=0$, $W=3q^{2e_0+5}-q^{2e_0+4}-q^{2e_0+5}=q^{2e_0+4}(2q-1)$;
if $k=1$, $W=q^{2e_0-1}+2q^{2e_0+5}-q^{2e_0+4}-q^{2e_0+5}
=q^{2e_0-1}(1-q^5+q^6)$.
\end{proof}

\begin{corollary}
Every negative coefficient of any single term of $N_2(c)$ arises from one of
the three sandwich polynomials $2q-1$, $1-q^5+q^6$, $1-q$, each contributing
exactly one negative coefficient per rank-$3$ profile $c'$
(namely $-q^{2e_0+4}(Q_{1,c'}+1)$ in total).
\end{corollary}

\section{Equidistribution at level $n$}

\begin{theorem}[Scalar reduction of divisibility]\label{thm:scalar}
Let $\gcd(d,3)=1$ and $\zeta$ be any root of $1+q^n+q^{2n}$, so that
$\omega_\zeta:=\zeta^n$ is a primitive cube root of unity. Since
$\EMD(c,c')\equiv\delta(c')-\delta(c)\pmod 3$ with $\delta(c)=(c_0-c_1)\bmod3$,
\[
N_n(c)\big|_{q=\zeta}
=\omega_\zeta^{-\delta(c)}\;S_n(\zeta),\qquad
S_n(\zeta):=\sum_{c'}\omega_\zeta^{\delta(c')}\,\mathrm{BR}_n(c')\big|_{q=\zeta}.
\]
Hence divisibility of all $|{\cal P}_d|$ numerators $N_n(c)$ by $1+q^n+q^{2n}$
is equivalent to the single scalar identity $S_n(\zeta)=0$ per root $\zeta$.
\end{theorem}

\begin{proof}
Immediate from $\zeta^{n\EMD}=\omega_\zeta^{\EMD}$ and the mod-$3$ formula for
$\EMD$ (Layer 1, Seed 7). The identity $S_n(\zeta)=0$ itself follows from
Welsh's polynomiality of $Q_n$; at $n=2$ it was also verified directly
($S_2(\omega)=S_2(-\omega)=0$ for $d=4,5,7$). This is the level-$n$
generalization of the EMD Equidistribution Theorem: the $c$-dependence factors
out completely, leaving one global charge-conservation identity.
\end{proof}

\section{What remains open}

\begin{conjecture}[Numerator positivity]\label{con:num}
$N_n(c)=(1+q^n+q^{2n})\,Q_{n,c}\ \ge\ 0$ coefficientwise for all $n,c$
(when $\gcd(d,3)=1$).
\end{conjecture}

Verified: $d=4$, $n\le5$, all $15$ profiles (exact arithmetic);
$d=5,7$, $n\le2$, all profiles. Two further facts constrain any proof:
(i) $(1+q+q^2)Q_2\ge0$ holds but $(1-q+q^2)Q_2\ge0$ \emph{fails} for every
profile tested, so positivity descends from $N_2$ through the $\Phi_6$
division first and the $\Phi_3$ division last;
(ii) the profile monotonicity $Q_{1,c'(J)}\ge q^a Q_{1,c'}$ is false for all
$a\ge1$, so no termwise pairing inside Theorem~\ref{thm:shape} can work:
positivity of $N_2$ is an irreducibly global, EMD-weighted phenomenon.
By Theorem~\ref{thm:shape} the required cancellation is now confined to a
single negative coefficient per rank-$3$ profile, a drastic reduction of the
search space for a sign-reversing injection.

\end{document}
